\scanpage{43}\mspage{42}
\uline{Cor. 2} Soit $F$ une figure. Deux lieux dans \deleted{\uline{Cor. 2} Soit $F$ une figure. Deux} l'ombre de $F$, qui appartiennent à des ombres ouvertes de strates distinctes, sont \uline{disjoints}.

Soient $\deleted{X}F$ une figure, \deleted{$X\in\unsure{\mathcal M|M}$ un raffinement de $F$,} \inserted{$F'$ un raffinement} de $F$. Je \unsure{veux|vaux} \unsure{definir|définir} une application
\[
 F'\longrightarrow F
\]
en associant à \unsure{tt|tout} $X'\in F'$ \unsure{\{l'unique\}|(l'unique)} $X\in F$ tel que
\[
 \widetilde{X'}^\circ\subset\widetilde X^\circ
\]
(ou encore le \uline{plus petit} \deleted{\unsure{des|x|X}} $\unsure{X'|X}\in F$ tel que $\widetilde{X'}^\circ\subset\widetilde X$). \unsure{Savons-nous} déjà qu'un tel $X$ existe, i.e. que \unsure{la|R} relation d'équiv. canonique sur l'ombre de $F'$, est plus fine que celle induite par l'ombre de $F$ ? Soit $X$ un élément de $F$ \uline{minimal}, parmi ceux satisfaisant $X'\ll X$ (existe, à cause de l'axiome de dim), \unsure{savons nous|savons-nous} que dès lors $\widetilde{X'}^\circ\subset\widetilde X^\circ$ ? \deleted{\unsure{Soit}} On sait que $\widetilde{X'}\subset\widetilde X$ donc si $x\in X'^\circ\inserted{\subset\widetilde{X'}}$, on a $x\in\widetilde{\deleted{\unsure{X|Y}}Y}^\circ\subset\widetilde Y$ pour $Y\unsure{\leq|\subseteq}X$, on veut prouver que $Y=X$. Il semblerait qu'il manque quelque chose.
